% GENERATED_BY: translate_oc_core_latex_targets.py
% AI_ASSISTED_TRANSLATION_DRAFT: TRUE
% THEOREM_REVIEW_STATUS: PENDING
% THEOREM_REVIEW_MODE: FORMAL_AI_GATE__LATEX_AWARE_MT
% THEOREM_REVIEW_REVIEWER_ID:
% THEOREM_REVIEW_AUTHORITY_CLASS:
% THEOREM_REVIEWED_AT:
% THEOREM_REVIEW_DOSSIER_REF:
% SOURCE_LANGUAGE: EN
% TARGET_LANGUAGE: DE
% SOURCE_EN_REF: content/cycles/cycles_k0.tex
% ================================================================
% ==== FILE: content/cycles/cycles_k0.tex
% ================================================================

\subsubsection{Zyklen auf \texorpdfstring{$K_0$}{K_0}}
\label{sec:cycles-k0}

Ebene $K_0$ stellt den metaontologischen Hintergrund dar
Ontology of Continua.
Es besitzt keine Geometrie, keine Zeit, keinen zulässigen Zustandsraum im
Sinn von $\Omega(K)$, keine Flüsse $J$, keine Schwellen $\Theta$ und nein
Operatoren der Evolution.
Aus diesem Grund kann der Begriff eines Zyklus auf $K_0$ nicht definiert werden.

\subsubsection{Abwesenheit von Zustandsraum und Zeit}

Zyklen erfordern:
\begin{enumerate}
  \item ein nicht leerer Zustandsraum $\Omega(K)$;
  \item eine Vorstellung von Zeit oder Parametrisierung $t\in[0,\tau]$;
  \item ein Flussfeld $J$, das geschlossene Trajektorien erzeugt.
\end{enumerate}

Keine dieser Strukturen existiert auf $K_0$.
Formell:
\[
  \Omega(K_0)=\emptyset, \qquad
  \partial\Omega(K_0) \text{ undefined}, \qquad
  J_{K_0} \text{ undefined},
\]
und es existiert kein Parameter $t$, der eine Flugbahn indizieren könnte.

Daher ein Kreislauf
\[
  C:[0,\tau]\to\Omega(K), \quad C(0)=C(\tau),
\]
lässt sich nicht formulieren.

\subsubsection{Struktureller Grund}

Ebene $K_0$ ist kein Kontinuum, sondern eine logische Bedingung der Möglichkeit
für Kontinua der Ebenen $K_1$ und höher.
Es legt keine dynamische oder geometrische Struktur fest, aus der Zyklen hervorgehen
entstehen könnte.
Alle zyklusbezogenen Begriffe – Länge, Stabilität, Abstand zur Grenze,
Zykluskomplexe oder Beiträge zur Kontinuität – sind streng genommen
nicht anwendbar.

\subsubsection{Auswirkungen auf höhere Ebenen}

Das Fehlen von Zyklen auf $K_0$ impliziert:
\begin{itemize}
  \item Der Zeitbegriff $\tau(K)$ beginnt erst bei $K_1$, wobei a
        Die kontinuierliche Achse wird zuerst definiert.
  \item die ersten nicht-trivialen Zyklen erscheinen bei $K_1$ als trivial oder
        entartete Schwingungsstrukturen eines eindimensionalen Feldes;
  \item Strukturzyklen im vollen Sinne (Stabilität, Abstand zu
        $\partial\Omega$, Interaktion mit Schwellenwerten) entstehen erst bei
        Ebenen $K_2$ und höher;
  \item Die Geburt von Zyklen ist selbst eine Signatur der Phase
        Übergang $K_0\to K_1$.
\end{itemize}

\subsubsection{Summary}

Auf $K_0$ sind keine Zyklen vorhanden.
Die Zyklusstruktur wird erst ab $K_1$ sinnvoll, wobei Zeit,
Zustandsraum und Strömungen erscheinen zuerst.
Dementsprechend zeichnet die Datei \texttt{cycles\_k0.tex} ein Formales auf
\emph{Nichtexistenz-Anweisung}, erforderlich für die interne Konsistenz
der Zyklustaxonomie über die gesamte Hierarchie der Kontinua hinweg.
